\documentclass[11pt]{article}

\usepackage[final]{acl}

\usepackage{times}
\usepackage[T1]{fontenc}
\usepackage[utf8]{inputenc}
\usepackage{microtype}
\usepackage{inconsolata}
\usepackage{booktabs}

\title{Creeping Crawler}

\author{Lapo Siciliani \\
  \texttt{2007890} \\\And
  Fabio Priori \\
  \texttt{1938446} \\\And
  Tamer Hayek \\
  \texttt{1897438} \\}

\begin{document}
\maketitle

% ─── 1. Objective ────────────────────────────────────────────────────────────

\section{Project Objective}

Creeping Crawler is an end-to-end pipeline for acquiring and analysing textual content from heterogeneous web sources, designed to feed RAG systems and language models. The system downloads web pages via Crawl4AI, transforms them into structured text through domain-specific parsers, and evaluates extraction quality by comparing the output against a manually built Gold Standard.

% ─── 2. Code Organisation ────────────────────────────────────────────────────

\section{Code Organisation}

The project is split into two Docker containers orchestrated by \texttt{docker-compose.yaml}: \textbf{backend} (FastAPI, port 8003) and \textbf{frontend} (FastAPI + Jinja2, port 8004).

\paragraph{Backend}

\begin{itemize}
  \item \textbf{\texttt{src/lib/crawling/}} — Crawl4AI-based acquisition module. \texttt{crawler.py} exposes three public functions: \texttt{fetch\_page} (live crawl), \texttt{fetch\_page\_from\_html} (parsing from already-available HTML via the \texttt{raw:} prefix), and \texttt{fetch\_page\_for\_url} (uses the HTML saved in the Gold Standard as the primary source; falls back to a live crawl only if unavailable). This design choice is motivated by the dynamic nature of the assigned domains: sites such as XE frequently update currency data, so re-crawling live would introduce drift relative to the manually built \texttt{gold\_text}, polluting the evaluation. The \texttt{domains/} sub-folder contains a dedicated \texttt{CrawlerRunConfig} for each domain, excluding irrelevant tags, CSS selectors, and forms; \texttt{CacheMode.BYPASS} ensures always-fresh content.

  \item \textbf{\texttt{src/lib/parsers/}} — Abstract \texttt{ContentParser} interface with domain-specific implementations in \texttt{domains/} (Wikipedia, ESPN, CNBC, XE). Each parser operates on the Markdown produced by Crawl4AI: it stops collection at boilerplate sections (e.g.\ ``Notes'', ``Related Content''), removes residual patterns via regular expressions (footnotes, navigation links, captions), and returns clean text in Markdown format.

  \item \textbf{\texttt{src/lib/evaluation/}} — Three-module evaluation pipeline: \texttt{tokens.py} (Markdown stripping via mistune/BeautifulSoup, Unicode normalisation), \texttt{token\_level.py} (set-based metrics: precision, recall, F1), \texttt{similarity.py} (vector metrics: cosine, Jaccard, excess ratio).

  \item \textbf{\texttt{src/lib/gold\_standard/}} — Access to \texttt{domains.json} and the JSON files in \texttt{gs\_data/}, one per domain, each containing representative entries with \texttt{url}, \texttt{domain}, \texttt{title}, \texttt{html\_text}, and \texttt{gold\_text}.

  \item \textbf{\texttt{src/routes/}} — Four FastAPI routers (\texttt{domains}, \texttt{parse}, \texttt{gold}, \texttt{evaluate}), each with dedicated Pydantic schemas in \texttt{src/schemas/}.
\end{itemize}

\paragraph{Frontend}

FastAPI server with Jinja2 rendering, two routes:

\begin{itemize}
  \item \textbf{\texttt{GET /}} — Home page: fetches the list of Gold Standard URLs from the backend via \texttt{/gold\_standard\_urls} and presents it as a dropdown menu grouped by domain (with \texttt{www.} stripped). The user can either select a URL from the menu or type an arbitrary one into the free-text field.

  \item \textbf{\texttt{GET /compare?url=\textless{}url\textgreater{}}} — Results page: calls \texttt{/parse} on the backend and produces clean text via \texttt{strip\_markdown}. If the URL is in the Gold Standard, it also fetches \texttt{gold\_text} and computes evaluation metrics. Displays four panels (raw HTML, parsed Markdown, clean text, Gold Standard) and, when available, an interactive Monaco Editor side-by-side diff between the extracted text and the Gold Standard.
\end{itemize}

% ─── 3. Parser Evaluation ───────────────────────────────────────────────────

\section{Parser Evaluation}

\paragraph{Methodology.} Each parser is evaluated on 5 URLs per domain via \texttt{GET /full\_gs\_eval}. Both the extracted text and the Gold Standard are passed through \texttt{strip\_markdown} before comparison, normalising typographic quotes, Unicode dashes, and HTML entities. Set-based metrics operate on sets of unique tokens (case-sensitive, space-separated); vector metrics (cosine, Jaccard, excess ratio) operate on frequency counts via \texttt{Counter}.

\paragraph{Results.} The table below reports average values across the entire Gold Standard for each domain (obtained via \texttt{GET /full\_gs\_eval?domain=<domain>}).

\begin{table}[h]
\centering
\small
\begin{tabular}{lcccccc}
\toprule
\textbf{Domain} & \textbf{P} & \textbf{R} & \textbf{F1} & \textbf{Cos} & \textbf{Jac} & \textbf{ExR}\\
\midrule
it.wikipedia.org  & -- & -- & -- & -- & -- & -- \\
www.espn.com      & -- & -- & -- & -- & -- & -- \\
www.cnbc.com      & -- & -- & -- & -- & -- & -- \\
www.xe.com        & -- & -- & -- & -- & -- & -- \\
\bottomrule
\end{tabular}
\caption{Average metrics per domain (P=Precision, R=Recall, F1, Cos=Cosine, Jac=Jaccard, ExR=Excess Ratio).}
\end{table}

\paragraph{Qualitative considerations.} The Crawl4AI configuration via \texttt{CrawlerRunConfig} (domain-specific CSS selector exclusion) is the first filtering layer. The Python parser then acts on the residual Markdown: for Wikipedia, service sections (``Notes'', ``Bibliography'', ``See also'') are cleanly cut off; for ESPN and CNBC, regex patterns remove video promotions, isolated links, and captions; for XE, the structured content is already nearly noise-free after configuration. The visual diff in the Monaco Diff Editor in the frontend guided iterative tuning of the exclusion patterns.

% ─── 4. Bonus/Extra ──────────────────────────────────────────────────────────

\section*{Bonus/Extra}

\paragraph{Similarity Eval (additional metrics).} Beyond the mandatory \texttt{token\_level\_eval}, a second group of metrics based on frequency vectors has been implemented: \emph{cosine similarity} (measures token distribution), \emph{Jaccard similarity} (penalises both extra and missing tokens), and \emph{excess ratio} (fraction of extracted tokens not covered by the Gold Standard — lower is better). These values complement the set-based evaluation by detecting redundant text not captured by precision/recall.

\paragraph{Script \texttt{crawl\_gs.py} and \texttt{Makefile}.} A standalone script downloads all Gold Standard URLs and saves the structured output in \texttt{gs\_results/} in five formats: raw and cleaned HTML, Markdown, stripped text, and parsed text. The \texttt{--update-json} option automatically updates the \texttt{html\_text} field in the Gold Standard JSON files to allow offline re-parsing without re-crawling. The Makefile automates the entire development cycle: conda environment creation, dependency installation (including Playwright/Chromium), server startup, crawling, requirements snapshot, and environment cleanup.

\paragraph{Monaco Diff Editor.} The frontend results page integrates Monaco Editor (the same engine behind VS Code) in side-by-side diff mode. When the URL is in the Gold Standard, the panel shows a token-by-token difference between the clean text extracted by the parser and the Gold Standard, highlighting exactly what was added or missed. This made parser tuning significantly faster compared to manual file-based comparison.

\paragraph{Additional endpoints.} Beyond the mandatory endpoints, \texttt{GET /gold\_standard\_urls} has been implemented; it returns the full list of URLs in the Gold Standard and is used by the frontend to populate the dropdown menu and determine whether a manually entered URL admits evaluation.

\end{document}
